% Propietario conservado: /Users/ruben/Documents/ChatGPT/jueces y controles/REVISION_ARGUMENTAL_APERTURAS_CIERRES_20260915/II/manuscript_es/sections/05_accion.tex
\section{Constantes de Planck et structure déterminantale de l'action}
\label{ii:sec:action}
\subsection{Échelle déterminantale d'action et valuation décimale}

L'échelle absolue et le correcteur relatif proviennent d'objets
distincts. La première utilise le bloc local du registre dodécaphasique
de pas trois,
\begin{equation}
 H_4=\begin{pmatrix}
 1&1&1&1\\1&1&-1&-1\\1&-1&1&-1\\1&-1&-1&1
 \end{pmatrix},\qquad
 G_H=\mathbb Z^4/H_4\mathbb Z^4.
 \label{ii:eq:el-h4}
\end{equation}
L'opération utilise un quotient local. Les trois blocs du registre
complet ne sont pas remplacés par une tensorisation qui multiplierait à
nouveau le nombre de feuillets de ce lecteur.

\begin{lemma}[Nombre de feuillets du quotient local]
La forme normale de Smith de \(H_4\) est
\(\operatorname{diag}(1,2,2,4)\). Par conséquent,
\[
 G_H\simeq\mathbb Z/2\mathbb Z\oplus\mathbb Z/2\mathbb Z
 \oplus\mathbb Z/4\mathbb Z,\qquad |G_H|=16.
\]
\end{lemma}
\begin{proof}
Le plus grand commun diviseur des entrées vaut un. Les mineurs d'ordre deux
sont pairs et l'un d'eux vaut \(-2\), de sorte que leur plus grand commun diviseur
vaut deux. Comme \(H_4H_4^{\mathsf T}=4I\) et \(\det H_4=-16\), les
cofacteurs valent \(\pm4\) ; le troisième diviseur déterminantal vaut quatre.
Le quatrième vaut \(16\). Les quotients successifs des diviseurs
\(1,2,4,16\) sont \(1,2,2,4\). Leur produit donne l'ordre du quotient.
\end{proof}

\begin{definition}[Lecteur déterminantal d'action]
La norme multiplicative de la section scalaire \(\alpha\), constante
sur les feuillets de \(G_H\), et une application de l'auto-échelle
\(\varphi\) produisent
\begin{equation}
 \Sigma_H=\varphi\,\mathrm N_{G_H}(\alpha),\qquad
 \mathrm N_{G_H}(\alpha)=\prod_{g\in G_H}\alpha=\alpha^{16}.
 \label{ii:eq:el-lector-determinantal}
\end{equation}
La carte décimale associée détermine
\(\operatorname{ord}_{10}(\Sigma_H)=\lfloor\log_{10}\Sigma_H\rfloor\).
\end{definition}

La puissance seize provient ainsi du nombre de feuillets du quotient, une
fois ce lecteur local fixé. Le facteur \(\varphi\) agit une fois sur
la base de la droite d'action. Le calcul de Smith, à lui seul, ne
sélectionnerait pas n'importe quelle fonctionnelle possible sur le quotient : la norme
multiplicative et l'auto-échelle font partie de l'opération déclarée.

\begin{theorem}[Décade de la section d'action]
Le lecteur \eqref{ii:eq:el-lector-determinantal}, appliqué aux coordonnées
HMT déjà engendrées, satisfait
\begin{equation}
 10^{-34}<\Sigma_H<10^{-33},\qquad
 \operatorname{ord}_{10}(\Sigma_H)=-34.
 \label{ii:eq:el-decada}
\end{equation}
\end{theorem}
\begin{proof}
Les intervalles rationnels suivants, beaucoup plus grossiers que les cylindres
disponibles, suffisent :
\[
 \frac{1618}{1000}<\varphi<\frac{1619}{1000},\qquad
 \frac{7297}{10^6}<\alpha<\frac{7298}{10^6}.
\]
La fonction \((x,a)\mapsto xa^{16}\) est croissante en ses deux arguments
positifs. Les deux inégalités entières
\[
 1618\cdot7297^{16}>10^{65},\qquad
 1619\cdot7298^{16}<10^{66}
\]
impliquent
\[
 10^{-34}
 <\frac{1618}{1000}\left(\frac{7297}{10^6}\right)^{16}
 <\Sigma_H
 <\frac{1619}{1000}\left(\frac{7298}{10^6}\right)^{16}
 <10^{-33}.
\]
La définition de l'ordre décimal conclut la preuve. Ni une
unité SI ni une masse de référence n'intervient dans ces inégalités.
\end{proof}

La section normalisée commence par
\(\Sigma_H=1.0462438381\ldots\times10^{-34}\). Sa mantisse ne
s'identifie pas à la coordonnée d'action avant ou après le retour.
Le lecteur déterminantal fixe l'ordre décimal ; les deux coordonnées de
retour s'obtiennent au moyen des lecteurs suivants.


\subsection{Sections d'action et quotient adimensionnel de retour}

Soit \(\mathcal L_S\) une droite orientée d'action et \(\mathcal U_S\)
sa base positive. Le caractère d'ordre cinq et le terme régional de
retour ont pour expressions
\begin{align}
 H_5={}&\frac12\sqrt{\frac{1000\alpha}{\varphi}}-\alpha
   +\frac9{16}\alpha^2-\frac59\alpha^3
   +\frac7{48}\alpha^4-\frac1{54}\alpha^5,
   \label{ii:eq:el-h5}\\
 D_A={}&\exp\!\left(-\frac{100\pi\alpha}{9}\right),
   \label{ii:eq:el-da}\\
 \mathcal C_\pi(\alpha)={}&169\alpha^6+
       \frac{D_A\alpha^7}{1-D_A\alpha},
   \label{ii:eq:el-cpi}\\
 \eta_{\mathrm{ret}}={}&H_5-\frac{90}{\pi}\mathcal C_\pi(\alpha).
   \label{ii:eq:el-eta}
\end{align}
Ce sont des lectures ultérieures de la fermeture dodécaphasique et de son
incidence régionale. L'entier \(169\) est conservé comme le sélecteur
régional déjà construit, et les coefficients de \(H_5\) comme ceux de son
caractère local ; ils ne sont pas remplacés par une interpolation de la constante de
Planck. Les formules explicitent exactement les lectures nécessaires
à l'étape électronique. Il n'est pas nécessaire d'introduire ici d'autres
fonctionnelles de la théorie des masses.

Pour \(0<\alpha<1\), on a \(0<D_A<1\), de sorte que
\(D_A\alpha<1\). Le terme rationnel de \eqref{ii:eq:el-cpi} conserve la
somme complète de la queue géométrique
\begin{equation}
 \frac{D_A\alpha^7}{1-D_A\alpha}
 =\sum_{n=0}^{\infty}D_A^{n+1}\alpha^{n+7}.
 \label{ii:eq:el-cola-regional}
\end{equation}
Il ne s'agit pas d'une troncature que l'on pourrait modifier silencieusement en changeant la
précision. Après le terme d'indice \(N\), la queue restante est
\(D_A^{N+2}\alpha^{N+8}/(1-D_A\alpha)\), ce qui permet de contrôler
l'évaluation sans recourir à une grandeur cible.

\begin{definition}[Sections d'action et lecteur de retour]
Les sections avant et après le retour sont
\begin{equation}
 \hbar_{\mathrm{pre}}=H_5\,10^{-34}\mathcal U_S,
 \qquad
 \hbar_{\mathrm{ret}}=\eta_{\mathrm{ret}}\,10^{-34}\mathcal U_S.
 \label{ii:eq:el-secciones-accion}
\end{equation}
Lorsque \(\eta_{\mathrm{ret}}>0\), on définit
\begin{equation}
 \delta_{\mathrm{ret}}=H_5-\eta_{\mathrm{ret}},\qquad
 R_{\mathrm{act}}=\frac{\hbar_{\mathrm{pre}}}{\hbar_{\mathrm{ret}}}
 =\frac{H_5}{\eta_{\mathrm{ret}}}.
 \label{ii:eq:el-ratio}
\end{equation}
\end{definition}

\begin{proposition}[Positivité et invariance du rapport]
Dans les intervalles utilisés dans la preuve de \eqref{ii:eq:el-decada},
\(0<\eta_{\mathrm{ret}}<H_5\), et donc
\(R_{\mathrm{act}}>1\). Le rapport ne change pas lorsqu'on remplace la base
\(\mathcal U_S\) par toute autre base positive commune.
\end{proposition}
\begin{proof}
Les termes de \(\mathcal C_\pi\) sont positifs. Pour obtenir une borne
élémentaire, il suffit d'utiliser \(\alpha<0.008\), \(\varphi<1.619\),
\(\alpha>0.007\), \(D_A<1\) et \(\pi>3\). La racine dans
\eqref{ii:eq:el-h5} est supérieure à un, tandis que la somme des termes
négatifs est inférieure à \(0.008001\) ; en particulier, \(H_5>0.99\).
De même,
\[
 0<\frac{90}{\pi}\mathcal C_\pi
 <30\left(169(0.008)^6+
           \frac{(0.008)^7}{1-0.008}\right)<10^{-6}.
\]
Ainsi, \(\eta_{\mathrm{ret}}>0.989999>0\) et
\(\eta_{\mathrm{ret}}<H_5\). Enfin, le changement de base divise
les deux coordonnées d'action par le même facteur positif, qui se simplifie dans
leur quotient.
\end{proof}

L'évaluation ultérieure fournit
\begin{align*}
 H_5&=1.0545718176461544696\ldots,\\
 \eta_{\mathrm{ret}}&=1.0545718169150390611\ldots,\\
 R_{\mathrm{act}}&=1.0000000006932817630\ldots.
\end{align*}
La différence et le quotient sont deux lectures des mêmes sections,
l'une additive et l'autre multiplicative. Elles ne représentent pas deux constantes de
Planck indépendantes. Le passage de l'action angulaire à l'action par tour
complet s'exprime, dans chaque section, par
\(h=2\pi\hbar\). La carte ultérieure
\(\mathcal U_S\mapsto\mathrm{J\,s}\) attribue une unité, mais ne
rétroagit pas sur l'ordre décimal déjà démontré.

Le retour régional peut aussi s'écrire dans une carte angulaire.
Cette reparamétrisation n'est pas une dépendance supplémentaire indispensable :
\eqref{ii:eq:el-h5}--\eqref{ii:eq:el-eta} fournissent directement
\(H_5\), \(\eta_{\mathrm{ret}}\) et \(R_{\mathrm{act}}\). Si l'on
introduit l'angle et le contre-angle, tous deux doivent être transportés dans la même
carte ; degrés et radians ne sont pas mélangés dans une différence.


% BEGIN OWNER /Users/ruben/Documents/ChatGPT/jueces y controles/REVISION_ARGUMENTAL_APERTURAS_CIERRES_20260915/II/manuscript_es/sections/revision_planck.tex
\subsubsection{Caractère d'action et provenance de ses coefficients}
\label{ii:sec:revision-planck}

La constante de Planck réduite, \(\hbar\), est l'objet physique qui
motive la réalisation de la section d'action angulaire. Sa position dans
la dépendance électronique doit être spécifiée à deux niveaux. La norme
déterminantale \(\Sigma_H\) détermine la décade de l'échelle. Le caractère
\(H_5\) et le retour régional déterminent ensuite les coordonnées des
sections d'action. Cette séparation permet d'exposer le résultat
nécessaire à l'électron sans anticiper les développements complets
d'autres secteurs de la Mécanique Dimensionnelle.

La provenance du caractère d'ordre cinq peut être explicitée au moyen
d'invariants déjà fixés dans le support. Dans la carte centrale, on considère
\[
 B_c=\begin{pmatrix}7&2\\2&7\end{pmatrix},
 \qquad \lambda_+=9,\quad\lambda_-=5.
\]
Le calendrier est de longueur \(12\), l'orbite du pas \(3\) est de
longueur \(4\), et le quotient de dénombrement pertinent est
\(104976/1944=54\). Les coefficients du caractère satisfont
\begin{equation}
 \frac9{16}=\frac{\lambda_+}{4^2},\quad
 \frac59=\frac{\lambda_-}{\lambda_+},\quad
 \frac7{48}=\frac{\operatorname{tr}(B_c)/2}{4\cdot12},\quad
 \frac1{54}=\frac{1944}{104976}.
 \label{ii:eq:revision-h5-coeficientes}
\end{equation}
L'attribution de ces invariants aux degrés \(2,3,4,5\), avec
l'alternance des signes de \eqref{ii:eq:el-h5}, appartient à la définition du
caractère analytique utilisé. L'expliciter est essentiel : l'ensemble des
cardinaux, considéré sans le caractère ni l'ordre de ses degrés,
admettrait d'autres expressions. Le contenu constructif réside dans la
composition du support, de l'orientation et de la règle analytique déclarée.

\subsubsection{Continuation régionale et équation de renouvellement}

Le sélecteur \(169\) détermine l'impulsion de degré six. Le prolongement
ultérieur s'organise selon le transfert \(D_A\) et une normalisation
unitaire de la droite déterminantale. Ces règles admettent une formulation
en séries formelles avant la substitution \(z=\alpha\), avec \(D_A\) déjà
construit. Écrivons
\[
 \mathscr C(z)=169z^6+R(z),\qquad R(z)\in z^7\mathbb R[[z]].
\]
La concaténation d'un prolongement augmente le degré d'une unité et
multiplie son poids par \(D_A\). Le premier prolongement strict a pour
poids \(D_A\). Son équation de renouvellement est donc
\begin{equation}
 R(z)=D_Az^7+D_AzR(z).
 \label{ii:eq:revision-renovacion}
\end{equation}

\begin{proposition}[Unicité formelle de la continuation normalisée]
L'équation \eqref{ii:eq:revision-renovacion} détermine une unique série et
\[
 \mathscr C(z)=169z^6+\frac{D_Az^7}{1-D_Az}
             =169z^6+\sum_{n=7}^{\infty}D_A^{n-6}z^n.
\]
La réalisation est analytique pour \(|D_Az|<1\).
\end{proposition}
\begin{proof}
L'élément \(1-D_Az\) a pour terme constant un et est inversible
dans l'anneau des séries formelles. Résoudre l'équation produit la série
affichée. De manière équivalente, le coefficient de degré sept vaut \(D_A\)
et chaque coefficient ultérieur vaut \(D_A\) fois le précédent. La série
géométrique converge absolument dans le disque indiqué.
\end{proof}

La normalisation du premier prolongement est une donnée opératoire
indispensable. Si l'on ne conservait que l'impulsion de degré six et
le rapport de concaténation, les séries
\[
 169z^6+\lambda\frac{D_Az^7}{1-D_Az}
\]
partageraient encore ces deux données pour tout \(\lambda\).
La règle unitaire fixe \(\lambda=1\) avant l'évaluation.
La contribution de chaque règle à l'unicité de l'opérateur de lecture
est ainsi localisée, et l'expression rationnelle conserve la continuation complète.

\subsubsection{Constante de Planck réduite et retour de la section d'action}

Dans la carte dimensionnelle de \eqref{ii:eq:el-secciones-accion}, la section
antérieure au retour est la construction du modèle destinée à la comparaison
avec la constante de Planck réduite :
\[
 \hbar_{\mathrm{pre}}=H_5\,10^{-34}\mathcal U_S.
\]
La section postérieure conserve également la compensation régionale
\(90\mathscr C(\alpha)/\pi\). Toutes deux appartiennent à la même droite
d'action et leur quotient \(R_{\mathrm{act}}\) mesure l'effet multiplicatif
du retour. Le passage à \(h\) par \(2\pi\) utilise la relation
entre la paramétrisation angulaire et un tour complet.

Cette construction intervient effectivement dans
\(R_{\mathrm{act}}^{80-54\alpha+6\Delta_4}\). L'électron conserve
sa coordonnée centrale et sa trajectoire ; l'échelle chronologique ne
remplace pas cette structure par un quantum indifférencié. On ne
multiplie pas non plus à nouveau par \(10^{-34}\) une masse déjà exprimée en MeV :
la décade appartient à la section d'action, se simplifie dans son quotient
sans dimension et ne réapparaît que là où agit la réalisation
dimensionnelle correspondante.

La comparaison numérique de \(H_5\), de la section retournée et de
\(h/(2\pi)\) est présentée dans la section~\ref{ii:sec:revision-planck-contraste}.
On distingue ainsi la dénomination de l'objet physique, l'équation que le
modèle lui attribue et la précision avec laquelle cette équation le reproduit.

% END OWNER

